\section{Arquitectura de solución}
\label{sec:arquitectura}
\textit{Responsable: Benjamín (diagrama: Alexander, DOC-01).}

Nubeduca corre en la región \texttt{us-east-1} sobre AWS Academy Learner Lab, usando solo el rol
\texttt{LabRole}. Cada microservicio tiene su propia API en Amazon API Gateway, sus funciones AWS
Lambda (Python 3.13) y su tabla Amazon DynamoDB. Los microservicios no leen la base de datos de otro:
se integran por eventos en Amazon EventBridge o, para comandos síncronos, por la API del servicio
dueño.

\evidencia{arquitectura.png}{Arquitectura de solución de Nubeduca. Fuente editable:
\texttt{nubeduca-arquitectura.drawio} en el repositorio \texttt{.github}.}{fig:arquitectura}

\subsection{Cómo fluye una petición}
\begin{enumerate}
  \item El navegador carga el frontend desde AWS Amplify Hosting.
  \item El usuario inicia sesión en Amazon Cognito; el ID token trae \texttt{custom:tenant\_id} y
        \texttt{custom:role}.
  \item Cada llamada va a la API Gateway del microservicio dueño; el authorizer de Cognito valida el
        token y la Lambda toma el tenant del token, nunca del cuerpo.
  \item La Lambda escribe el cambio en su tabla y, en la misma transacción, el evento (ítem
        \texttt{EVENTO\#}) validado contra su esquema. Los archivos suben directo a Amazon S3 con URL
        prefirmada.
  \item DynamoDB Streams y EventBridge Pipes llevan cada evento al bus \texttt{lms-<stage>}; el Archive
        permite reprocesar.
  \item Una regla por consumidor copia el evento a su cola Amazon SQS con DLQ; la Lambda consumidora lo
        procesa una sola vez (idempotente).
  \item ms-calificaciones publica además las notas en Amazon SNS FIFO, que llegan en orden a la cola
        SQS FIFO de ms-notificaciones.
  \item ms-notificaciones avisa por correo con Amazon SNS.
  \item Amazon CloudWatch y AWS X-Ray observan todo; Serverless Framework despliega cada stack.
\end{enumerate}

\subsection{Principios de diseño}
\begin{enumerate}
  \item Cada servicio es dueño de sus datos.
  \item Los eventos son hechos en pasado (\texttt{submission.graded}); los comandos son llamadas
        síncronas al servicio dueño.
  \item Coreografía por defecto: los servicios reaccionan a eventos sin conocerse entre sí.
  \item \textit{Transactional outbox}: el cambio y el evento se guardan en la misma transacción.
  \item Consumidores idempotentes: la entrega es \textit{at-least-once}.
  \item Orden solo donde importa (notas, con SNS/SQS FIFO).
  \item Contratos versionados: un cambio incompatible es un tipo nuevo (\texttt{.v2}).
  \item Multi-tenant desde el día uno: \texttt{tenant\_id} en el token, en cada evento y en la clave
        de partición.
\end{enumerate}

\subsection{Decisiones de arquitectura (ADR)}
\begin{table}[htbp]
\centering\small
\begin{tabularx}{\textwidth}{lX}
\toprule
ADR & Decisión \\
\midrule
ADR-001 & 6 microservicios desplegables, cada uno agrupa varios dominios lógicos. \\
ADR-002 & Construir desde cero (no envolver Moodle). \\
ADR-003 & EventBridge + SQS/SNS como broker de eventos. \\
ADR-005 & Modelo \textit{pool}: los 3 tenants comparten tablas, con \texttt{tenant\_id} en la clave. \\
ADR-006 & Lambda + API Gateway para el cómputo. \\
ADR-007 & DynamoDB en los 6 microservicios (sin VPC ni RDS). \\
ADR-008 & Outbox con EventBridge Pipes en staging y prod; la Lambda relay queda como vuelta atrás. \\
\bottomrule
\end{tabularx}
\caption{Decisiones de arquitectura registradas en el repositorio \texttt{.github}.}
\label{tab:adr}
\end{table}
